\origpage{53}{GRAPH MINORS. IX}
\textit{Proof.} Let $\mathcal{Q} = \{Q_1, \dots, Q_{p+4}\}$, and choose a crooked transaction $\mathcal{P}$ in $(G, \Omega)$ such that the statement $S$ of the previous section is true. By (3.4), $|\mathcal{P}| = p$, and by (4.6), $Q_j \cap \bigcup \mathcal{P}$ is null for some $j$ ($1 \le j \le p+4$). We assume that (ii) does not hold, and so $V(\mathcal{P}) \cap \Omega(a_j, b_j) \ne \emptyset$. Hence by (5.1) there is a cross $\{P_1, P_2\} \subseteq \mathcal{P}$, with all four ends in $\Omega(a_j, b_j)$. Since $Q_j$ is separating it follows that $P_1 \cap G_2$ and $P_2 \cap G_2$ are both null, and (i) holds, as required. $\square$

\stmt
(5.3) \textit{Let $Q$ be an induced $\Omega$-path in $(K, \Omega)$ with ends $a, b$ such that $\Omega(a, b) = \emptyset$ and $\bar{\Omega} - \{a, b\}$ catches $K \setminus V(Q)$, and no internal vertex of $Q$ is in $\bar{\Omega}$. Let $r \in V(Q) - \{a, b\}$, and let $\Omega_1$ be the cyclic permutation in $V(K)$ with $\bar{\Omega}_1 = \bar{\Omega} \cup \{r\}$ and $\Omega_1 \mid \bar{\Omega} = \Omega$, and with $(a, r, b)$ $\Omega_1$-clockwise. If $(K, \Omega_1)$ has a cross then $(K, \Omega)$ has a cross or tripod.}

\textit{Proof.} Let $\{P_1, P_2\}$ be a cross in $(K, \Omega_1)$, chosen with $P_1 \cup P_2 \cup Q$ minimal. Let $P_i$ have ends $s_i, t_i$ ($i = 1, 2$). We may assume that $\{P_1, P_2\}$ is not a cross of $(K, \Omega)$, and so $r \in \{s_1, t_1, s_2, t_2\}$.

\hspace*{12pt}(1) \textit{Either $(K, \Omega)$ has a cross or $a, b \in \{s_1, t_1, s_2, t_2\}$.}

For certainly $Q$ meets $P_1 \cup P_2$; let $A$ be a minimal subpath of $Q$ from $a$ to $V(P_1 \cup P_2)$, with ends $a, a'$ where $a' \in V(P_1)$, say. Let $S, T$ be the two subpaths of $P_1$ from $a'$ to $s_1, t_1$, respectively. One of $\{A \cup S, P_2\}$, $\{A \cup T, P_2\}$ is a cross in $(K, \Omega_1)$ and we assume the former without loss of generality. From the minimality of $P_1 \cup P_2 \cup Q$ we deduce that $T \subseteq Q$ and in particular, $t_1 \in V(Q) \cap \bar{\Omega}_1 = \{a, b, r\}$. We may assume that $(A \cup S, P_2)$ is not a cross in $(K, \Omega)$ and so $r$ is an end of one of $A \cup S, P_2$; and it follows that $r \notin V(T)$, since $P_1 \ne T$. Since $t_1 \in \{a, b, r\}$ and $r \notin V(T)$ we deduce that $a = t_1 \in \{s_1, t_1, s_2, t_2\}$, and similarly $b \in \{s_1, t_1, s_2, t_2\}$, as claimed.

Since $r \in \{s_1, t_1, s_2, t_2\}$, we may assume that $r = s_2$ without loss of generality. Then $t_2 \in \Omega_1(b, a)$, since $\Omega_1(a, b) = \{r\}$ and $P_1$ crosses $P_2$. From (1), we may assume that $s_1 = a$ and $t_1 = b$. Let $R_2$ be the minimal subpath of $P_2$ between $\{t_2\}$ and $V(Q)$, and let $R_2$ have ends $t_2, c$, say, where $c \in V(Q)$. Let $C$ be the subpath of $Q$ between $a$ and $c$, and let $D$ be the subpath between $c$ and $b$. Now $a, b \in V(P_1)$, and so $P_1$ meets both $C$ and $D$. Choose a minimal subpath $R_1$ of $P_1$ between $V(C)$ and $V(D)$. Now $R_1$ has more than one edge, since $Q$ is induced, and so $V(R_1) \nsubseteq V(Q)$. But $\bar{\Omega} - \{a, b\}$ catches $K \setminus V(Q)$, and so there is a $V(R_1) \to (\bar{\Omega} - \{a, b\})$ path $P \subseteq K$ which avoids $V(Q)$. If $P \cap R_2$ is null then $(K, \Omega)$ has a cross $\{M_1, M_2\}$ where $M_1$ is the union of $P$, part of $R_1$ and part of $Q$, and $M_2$ is the union of $R_2$ and part of $Q$. If $P \cap R_2$ is non-null then $(K, \Omega)$ has a tripod $\{N_1, N_2, N_3\}$ the ``ends'' of which are $c$ and the vertex of $P \cap R_1$, where $N_1, N_2 \subseteq Q \cup R_1$ and $N_3 \subseteq P \cup R_2$. This completes the proof. $\square$

We shall also need the following extension of (2.3).
